% !TeX root = ../tfg.tex
% !TeX encoding = utf8
%
%*******************************************************
% Capítulo 3. Aprendizaje supervisado y redes profundas
%*******************************************************

\chapter{Aprendizaje supervisado y redes profundas}\label{cap:aprendizaje}

% Capítulo de fundamentos: todo es esqueleto. Notación a fijar: espacio de entrada
% \mathcal{X} \subseteq \R^{H \times W \times 3}, etiquetas \mathcal{Y} = \{1, \dots, C\} con C = 14,
% distribución conjunta P(X, Y), clasificador f_\theta: \mathcal{X} \to \R^C (logits),
% extractor de características h_\theta: \mathcal{X} \to \R^D (D = 768 en el Swin-S).

\section{El problema de clasificación supervisada}\label{sec:clasificacion}

\begin{definicion}[Problema de clasificación]\label{def:clasificacion}
% Muestra i.i.d. \{(x_i, y_i)\}_{i=1}^n de P(X, Y); clasificador g(x) = \argmax_k f_\theta(x)_k;
% riesgo esperado con pérdida 0-1 y riesgo empírico.
\pendiente{definición formal}
\end{definicion}

% - Clasificador de Bayes y su optimalidad respecto a la pérdida 0-1.
\begin{proposicion}[Optimalidad del clasificador de Bayes]\label{prop:bayes}
\pendiente{enunciado}
\end{proposicion}
\begin{demostracion}
\pendiente{demostración}
\end{demostracion}

\section{Función de pérdida: entropía cruzada y \emph{label smoothing}}\label{sec:perdida}

\begin{definicion}[Función softmax]\label{def:softmax}
% \softmax(z)_k = e^{z_k} / \sum_j e^{z_j}.
\pendiente{definición}
\end{definicion}

\begin{definicion}[Entropía cruzada]\label{def:entropia-cruzada}
\pendiente{definición}
\end{definicion}

% - Interpretación como estimación de máxima verosimilitud.
\begin{proposicion}\label{prop:ce-mv}
% Minimizar la entropía cruzada empírica equivale a maximizar la log-verosimilitud del modelo
% categórico.
\pendiente{enunciado}
\end{proposicion}
\begin{demostracion}
\pendiente{demostración}
\end{demostracion}

\begin{definicion}[\emph{Label smoothing}]\label{def:label-smoothing}
% Pérdida (1-\varepsilon) NLL + (\varepsilon/K) \sum_j -\log \pi_j, con \varepsilon = 0,1 (valor por
% defecto de LabelSmoothingCrossEntropy de fastai, usado en el entrenamiento).
\pendiente{definición}
\end{definicion}

\begin{proposicion}[Óptimo de la pérdida con \emph{label smoothing}]\label{prop:ls-optimo}
% El óptimo asigna \pi = \varepsilon/K a las clases no verdaderas y 1 - \varepsilon + \varepsilon/K a la
% verdadera (0,00714 y 0,907 con \varepsilon = 0,1 y K = 14). Se usa en el capítulo 12 para explicar
% el «suelo» del softmax (results_conforme/mondrian/label_smoothing.json).
\pendiente{enunciado}
\end{proposicion}
\begin{demostracion}
% Minimización con multiplicadores de Lagrange sobre el símplex.
\pendiente{demostración}
\end{demostracion}

\section{Optimización: Adam y AdamW}\label{sec:adamw}

% - Descenso de gradiente estocástico por minilotes.
% - Adam (Kingma y Ba, 2014): estimaciones de primer y segundo momento con corrección de sesgo.
% - AdamW (Loshchilov y Hutter, 2019): decaimiento de pesos desacoplado; diferencia con la
%   regularización L2 en métodos adaptativos.
\begin{definicion}[Adam]\label{def:adam}
\pendiente{reglas de actualización}
\end{definicion}

\begin{definicion}[AdamW]\label{def:adamw}
\pendiente{reglas de actualización}
\end{definicion}

\begin{observacion}
% Con Adam, la regularización L2 no equivale al decaimiento de pesos (Loshchilov y Hutter, 2019).
\pendiente{observación}
\end{observacion}

\section{Arquitecturas}\label{sec:arquitecturas}

\subsection{Redes neuronales convolucionales}\label{sec:cnn}

% - Operación de convolución discreta, equivarianza a traslaciones, agrupación (pooling).
% - ResNet (He et al., 2016): conexiones residuales. EfficientNetV2 (Tan y Le, 2021).
\begin{definicion}[Convolución discreta]\label{def:convolucion}
\pendiente{definición}
\end{definicion}

\subsection{Transformers y Vision Transformer}\label{sec:transformer}

% - Atención por producto escalar: Attention(Q, K, V) = softmax(QK^T / \sqrt{d}) V.
% - Atención multicabeza, LayerNorm, MLP; ViT (Dosovitskiy et al., 2021): parches de 16 px.
\begin{definicion}[Atención por producto escalar]\label{def:atencion}
\pendiente{definición}
\end{definicion}

\subsection{Swin Transformer}\label{sec:swin}

% - Liu et al. (2021): atención en ventanas locales (7x7 parches), ventanas desplazadas entre
%   bloques consecutivos, fusión de parches (jerarquía), coste lineal en el número de parches.
% - Swin-S (swin_small_patch4_window7_224): características de 768 dimensiones tras el pooling,
%   con LayerNorm (relevante para ReAct, ASH y DICE en el capítulo 10).
\begin{proposicion}[Coste de la atención por ventanas]\label{prop:swin-coste}
\pendiente{enunciado: complejidad global frente a por ventanas}
\end{proposicion}

\subsection{ConvNeXt}\label{sec:convnext}

% - Liu et al. (2022): red convolucional modernizada con decisiones de diseño de los transformers.
\pendiente{ConvNeXt}

\endinput
